\section*{Regulamin ćwiczeń}

Zaliczenie ćwiczeń uzyskuje się na podstawie frekwencji oraz ocen ze sprawdzianów.

W semestrze odbywają się trzy sprawdziany. Każdy ze sprawdzianów ma miejsce po zakończeniu jednego z bloków tematycznych (Statyka, Kinematyka, Dynamika). Sprawdziany odbywają się w trakcie zajęć, na pierwszej godzinie ćwiczeń.

Warunkiem zaliczenia ćwiczeń jest uzyskanie pozytywnej oceny (min. 50\%) z każdego ze sprawdzianów oraz utrzymanie frekwencji na poziomie min. 80\%.

Za każde z zadań na sprawdzianie można uzyskać ocenę w skali 0-5:
\begin{itemize}
    \item 5 za rozwiązanie bezbłędne i kompletne,
    \item 2-4 za rozwiązanie kompletne ale z błędami (gradacja w zależności od liczby i wagi błedów),
    \item 0 za rozwiązanie niekompletne lub całkowicie błędne.
\end{itemize}

Ocena ze sprawdzianu to średnia arytmetyczna ze wszystkich zadań na danym sprawdzianie. Ocena z ćwiczeń to średnia arytmetyczna ze wszystkich sprawdzianów.

Pod koniec semestru planowana jest jedna zbiorcza poprawa wszystkich sprawdzianów. Do poprawy może przystąpić każdy. Nie trzeba pisać wszystkich trzech sprawdzianów, a jedynie te wybrane przez siebie. Czas poprawy to 90 minut, niezależnie od liczby pisanych sprawdzianów. Obowiązuje ostatnia uzyskana ocena.

Na każdym sprawdzianie można mieć ze sobą kalkulator oraz jedną kartkę A4 dwustronnie i własnoręcznie zapisaną dowolnymi wzorami i notatkami. Kartki wydrukowane nie będą dopuszczone. Korzystanie z telefonów jest zabronione.

Prowadzący ma prawo do wezwania wskazanych studentów do ustnej obrony swoich pisemnych prac.

Podstawowy podręcznik do ćwiczeń:
E Wittbrodt, S Sawiak, Mechanika ogólna, Wydawnictwo PG